% ===== 底本 概论（Overview，无章号）；片段 c01；源块 B001--B011 =====
% 译者：Lead。底本 Overview 用 h2（无 chapter 计数），故本译本用 \section*（不编号、进目录）。

\section*{概论}
\addcontentsline{toc}{section}{概论}

% [源 B001]
在 Kontsevich 关于同调镜像对称（Homological Mirror Symmetry，HMS）的著名猜想 [Kon94] 提出之后，
「如何对一个给定簇几何地构造其镜像」便成为数学家关心的重要问题。Strominger--Yau--Zaslow
提出了一个重要方案 [SYZ96]：给定带有拉格朗日环面纤维化 $X\to B$ 的 Kähler 流形 $X$，
其镜像应当是同一底 $B$ 上的 Kähler 流形 $\check{X}\to B$，同样带有拉格朗日环面纤维化，
并且「纤维互为对偶环面」。

% [Tutor Note: B001 的 SYZ 原始表述只要求「同一底 $B$ 上的对偶环面纤维」，并未要求 $B$ 带整仿射结构；
%   本讲义第 2 章补出了 $B^{reg}$ 上的整仿射结构，阅读时不要把「底 $B$」误当作普通光滑流形。]

% [源 B002]
这个方案首先在椭圆曲线上获得成功，相应的同调镜像对称（部分地）由 Polishchuk--Zaslow 验证 [PZ98]。
但人们很快发现：对一般的 Kähler 流形 $X$，找到光滑的拉格朗日环面纤维化几乎是不可能的——
可能出现\textbf{环面奇点}（toric singularities），使得纤维的维数降低；也可能出现奇异纤维。
Cho 与 Oh [CO06] 研究了环面奇点对 Fano 环簇的影响，证明环面奇点的存在会给出镜像的
Landau--Ginzburg 超势 $W$，其构造方式是数全纯圆盘与环面奇点的相交数。
对于容许只带环面奇点的典范拉格朗日环面纤维化的环簇，相应的同调镜像对称陈述先由
Fang--Liu--Treumann--Zaslow [FLTZ11] 对光滑射影环簇证明，随后由 Kuwagaki [Kuw20] 推广到所有环叠。

% [Tutor Note: B002 中「intersection numbers of holomorphic disks with the toric singularities」
%   指的是圆盘与环面边界除子的相交数（$\mu(u)=2\sum_i(u\cdot D_i)$，见第 3 章），
%   不是与奇点集本身的相交；翻译时保留「环面奇点」是为了照顾原句，但技术上应理解为与边界除子的相交数。]

% [源 B003]
奇异纤维中最简单、也最一般的情形是非退化奇点，它类似于 Lefschetz 纤维化的奇异纤维，
由把余维 $1$ 的子环面收缩成一点得到。Auroux [Aur07] 研究了 $X=\mathbb{C}^2\setminus\{xy=1\}$
这一情形，其拉格朗日环面纤维化是由 Lefschetz 纤维化 $\pi(x,y)=xy$ 构造出来的。

% [源 B004]
在这一情形下可以清楚地看到\textbf{墙穿越现象}（wall-crossing）：环面纤维化 $X\to\mathbb{R}^2$ 只有一个奇异值，
来自 Lefschetz 纤维化的奇点集 $\{xy=0\}\subseteq\mathbb{C}^2$，它对应于墙 $\{x=0\}\times\mathbb{R}\subseteq\mathbb{R}^2$。
Auroux [Aur07] 证明：对任意 $b=(b_1,b_2)$，只要 $b_1>1$ 或 $b_1<1$，用来计算
$X\subseteq\mathbb{P}^2$ 的镜像超势的圆盘计数就不同；于是两个拉格朗日环面的模空间以非平凡的
过渡函数粘合在一起。

% [Tutor Note: B004 原文作「$b_1>1$ or $b_1<1$」，字面上遗漏了 $b_1=1$——恰是墙本身所在处，
%   该处计数不收敛/需另作处理。本译本照译，不擅改；此处属原文表述不严谨而非印刷错误。]

% [源 B005]
尽管这一现象看起来神秘，Fukaya 在 [Fuk09] 中给出了一个可能的解释：把上述图像加细，
考虑奇异拉格朗日环面 $L$ 的 Floer 形变空间，其中局部系统扮演形变 $CF^{\ast}(L,L)$ 上
Floer $A_{\infty}$-结构的角色。对奇异拉格朗日子流形 $L\subseteq X$，$CF^{\ast}(L,L)$ 上的
$A_{\infty}$-结构通常是\textbf{有曲率的}，即有曲率项 $\mu^0$ 阻碍恒等式 $(\mu^1)^2=0$ 成立。
但可能存在元素 $b\in H^1(L,\Lambda_0)$（其中 $\Lambda_0$ 是\textbf{诺维科夫环}），使形变 $\mu_b$ 满足
$(\mu_b^1)^2=0$。这样的 $b$ 必须满足 \textbf{Maurer--Cartan 方程}
\[
  \sum_k\mathfrak{m}^k(b,b,\dotsb,b)=0,
\]
于是我们称满足此方程的元素为 \textbf{Maurer--Cartan 元素}，并用 $\mathcal{MC}(L)$ 记全体
Maurer--Cartan 元素构成的空间（等价地，$L$ 的无穷小形变空间）。

% [Tutor Note: B005 底本源文此处写 $\lambda_0$（小写）却指「诺维科夫环 $\Lambda_0$」，前后不一致；
%   本译本按语义排作 $\Lambda_0$，属底本笔误（已登记 README「原文笔误登记表」）。
%   另：同一段把 Floer 微分记作 $\mu^k$、把 Maurer--Cartan 方程里的乘法记作 $\mathfrak{m}^k$，
%   这是底本的记号混用；读者应把 $\mathfrak{m}^k$ 理解为 $A_\infty$-运算。]

% [源 B006]
对任一 Maurer--Cartan 元素 $b\in\mathcal{MC}(L)$，对 $(L,b)$ 定义了 $X$ 的 Fukaya 范畴中的一个对象。
Fukaya 提议：$X$ 的镜像应当是一个并集的商
\[
  \mathcal{MC}(\mathfrak{L}):=\bigcup_{u\in B}\mathcal{MC}(L_u)\big/H^1(L_u;2\pi\sqrt{-1}\mathbb{Z}),
\]
其中 $\mathfrak{L}:=\{L_u\}_{u\in B}$ 是给定的一族拉格朗日子流形，$B$ 是某个光滑流形，
其维数与 $H^1(L;\mathbb{R})$ 相同。进而，对任意给定对象 $(L',b')\in\mathcal{F}(X)$，
Yoneda 函子
\[
  (u,b)\mapsto CF^{\ast}((L',b'),(L_u,b);\Lambda)
\]
应当在 $\mathcal{MC}(\mathfrak{L})$ 上定义一个凝聚层复形。

% [Tutor Note: B006 的商 $\big/H^1(L_u;2\pi\sqrt{-1}\mathbb{Z})$ 是「按局部系统的单值化作用取商」，
%   若忘记这一步，$\mathcal{MC}(\mathfrak{L})$ 会多出整格的冗余；这是 SYZ 镜像构造里最常被漏掉的细节。]

% [源 B007]
本次学习讨论班的主要目标，是了解 SYZ 镜像构造当前的研究状况，尤其是 Abouzaid [Abo21]
关于光滑拉格朗日环面纤维化（从而对所有阿贝尔簇）的 SYZ 与同调镜像对称的工作，
以及 Abouzaid--Auroux--Katzarkov [AEK16] 关于沿热带极限附近超曲面爆破环簇的工作。

% [源 B008]
\subsection*{讨论班日程}
\addcontentsline{toc}{subsection}{讨论班日程}

% [源 B009]
讨论班在每周一 16:00--17:30 于 \textbf{N0.008} 举行，可与周二 13:00--14:00 的辛几何讨论班协调安排。
日程如下：

% [Tutor Note: B009 的「16-17:30」按底本照译；表中日期形如 20.10.25 即 2025 年 10 月 20 日（周一），
%   与「每周一」自洽，可据此核对后续日期与报告主题的对应关系。]

% [源 B010]
\begin{table}[htbp]
\centering
\footnotesize
\setlength{\tabcolsep}{3pt}
\begin{tabular}{@{}lp{4.7cm}lp{4.0cm}l@{}}
\toprule
日期 & 内容 & 报告人 & 参考文献 & 备注 \\
\midrule
20.10.25 & 引言 & Si-Yang Liu & & \\
27.10.25 & 拉格朗日环面纤维化 & Yunpeng Niu & [Eva23，第 1、2 章] & \\
03.11.25 & 由环面边界得到的 Landau--Ginzburg 超势 & Si-Yang Liu & [CO06] & \\
17.11.25 & 焦点-焦点奇点 & Yunpeng Niu & [Eva23，第 6 章]、[Sym03] & \\
24.11.25 & 由环面边界得到的 Landau--Ginzburg 超势（续） & Kyoungmo Kim & [CO06] & \\
01.12.25 & 近环面纤维化的 SYZ 镜像对称与 Auroux 系统 & Si-Yang Liu & [Aur07] & \\
26.01.26 & 族 Floer 上同调 & Maximilian Schimpf & [Abo14] & \\
02.02.26 & 光滑拉格朗日环面纤维化的镜像对称 & Maximilian Schimpf & [Abo21] & \\
\bottomrule
\end{tabular}
\tabtitle{1}{学习讨论班日程（2025--2026 学年）}
\end{table}

% [Tutor Note: B010 表格的「内容」列中，底本把 24.11.25 一场的 Landau--Ginzburg 拼作
%   「Landau-Ginburg」（少一个 z），本译本按通行拼写排印并登记为底本笔误；
%   另外底本此表后面还有两行（26.01.26、02.02.26 的补充报告）被 HTML 注释掉了，故未列入。]

% [源 B011]
% 底本此处为 HTML 注释中的两行日程（<!-- <tr>… 26.01.26 / 02.02.26 … -->），
% 在网页上不显示，且与本文件 B010 表中的行重复；为忠实起见不作排版，详见 README「已知偏差与限制」。
